%% FORMATO ARTÍCULO ELECTIVA EN MATEMATICAS - UNIVERSIDAD DE CUNDINAMARCA %%

% En el directorio de trabajo deben estar los archivos: estilo (tipo LaTeX Class), articulo (tipo LaTeX Class), éstos archivos no se deben modificar,  Bibfile (tipo BibTex Database) en donde el autor debe introducir los datos de las referencias en formato BibTex y  Formato.tex (éste archivo en donde se edita el contenido del documento


%-----------------------------------------------------
%PAQUETES
\documentclass{estilos}
\usepackage[pdftex]{graphicx}
\usepackage{verbatim}
\usepackage[centerlast,sc,footnotesize]{caption}
\usepackage{subfigure}
\usepackage{amsmath}
\usepackage[spanish, es-tabla]{babel} %Idioma español
\usepackage[utf8]{inputenc} %Acentos español
%--------------------------------------------------

%%%%%%%%%%% TITULO DEL ARTICULO (Los títulos, resúmenes y palabras clave se organizan dependiendo del idioma en el que se escriba el artículo. Entonces, si el artículo se escribe en español se debe colocar primero el título en español y luego en inglés. Si el artículo se escribe en inglés se coloca primero el título en inglés y luego en español. Así mismo para el resumen y palabras clave)
\title{TÍTULO DEL ARTÍCULO}

%%%%%%%%%%% AUTORES DEL ARTICULO
\author{Autor I\suprm1, Autor II \suprm2 }
% Por cada autor \suprm# es para referenciar la afiliacion del autor con el número #. Se debe incluir de cada autor el primer nombre, inicial del segundo nombre si lo tiene y el primer apellido. 

%%%%%%%%%%% INSTITUCION DE LOS AUTORES
\newcommand{\afiliacion}{\suprm 1 Afiliación autor I \\ \suprm 2 afiliación autor II }
% Escribir afiliación (Grupo de investigación, Dependencia, Universidad, País) de los autores dentro del corchete derecho. 
% Recuerde que el caracter \\ es para comenzar una nueva linea

%%%%%%%%%%% ENCABEZADOS
\newcommand{\autor}{Autor I y Autor II}
\newcommand{\tema}{Título corto del artículo}
% Autor 1 y Autor 2         -----> aparecen en el encabezado de las paginas pares
% Titulo corto del articulo -----> aparece en el encabezado de las paginas impares


%%%%%%%%%%% CORREOS ELECTRONICOS DE LOS AUTORES
\authorinfo{e-mail de contacto}
%A1 y A2 son las iniciales del nombre de los autores correspondientes

%%%%%%%%%%% PAGINA EN QUE COMIENZA EL ARTICULO (RESERVADO PARA EL EDITOR)
\setcounter{page}{1}

%%%%%%%%%% Estilo de la bibliografia segun la American Physics Society 
\bibliographystyle{apsrev4-1}
\begin{document}
  \maketitle

%%%%%%%%%%% RESUMEN DE SU ARTICULO EN ESPAÑOL
\begin{resumen}
\noindent Escribir resumen en español.
\end{resumen}

%%%%%%%%%%% PALABRAS CLAVE DE SU ARTICULO EN ESPAÑOL
\keywords{Palabras clave en español.}

%%%%%%%%%%% CUERPO DEL DOCUMENTO
% El texto debe dividirse en secciones, cada una con un encabezado (por ejemplo: Introducción; Parte Experimental, Materiales y Métodos o Desarrollo teórico; Resultados; Discusión; Conclusiones). Se recomienda que estas secciones sean breves y equilibradas.
\section*{Introducción}
En esta sección debemos orientar al lector hacia el tema de estudio y la
motivación por hacerlo elegido. Es decir, justificarlo o responder la pregunta
¿Por qué se realizó? Asimismo, debemos enunciar claramente el
propósito u objetivo del experimento.

\section*{Marco teórico}
Es aconsejable que incluyamos un marco teórico-experimental del tema
que estudiamos. Las referencias citadas aquí son ejemplos de cómo se debe referenciar un artículo \cite{1}, un libro \cite{2}, una conferencia \cite{3}, una página web \cite{4} y memorias de congreso \cite{5}.\\

Además, para unidades, abreviaturas y símbolos, se usará el Sistema Internacional de Unidades (m, kg, s, K), empleando sólo términos aceptados generalmente. Es necesario explicar las abreviaturas desconocidas cuando se usen por primera vez. Se debe poner especial cuidado al escribir los símbolos para que sean identificados claramente. En casos especiales, se especificará directamente con el autor el uso de fórmulas, caracteres especiales u otros.

Las ecuaciones deben estar debidamente numeradas y referenciadas en el cuerpo del texto como en la ec. (\ref{miecuacion}).

\begin{equation}
\label{miecuacion}
x = \frac{1- e^{t/\tau}}{\sqrt{mc^2/E_0}}\cos(\omega t-\phi)
\end{equation}

\section*{Metodología}
Aquí se describe los procedimientos seguidos.
Si el estudio tiene recolección de datos experimentales es importante incluir el 
instrumental usado. Para ésto, es útil incluir un esquema del diseño experimental
elegido (ver fig. \ref{mifigura1}). Además puede recurrirse a diagramas esquemáticos 
que muestren las características más importantes del arreglo experimental 
y la disposición relativa de los instrumentos. Indicar también cuáles variables
se miden directamente, cuáles se obtienen indirectamente y a cuáles tomamos como datos de
otras fuentes (parámetros físicos, constantes, etc.). También es aconsejable
describir las virtudes y limitaciones del diseño experimental.

\begin{figure}[H]
\begin{center}
\includegraphics[width=7cm]{figura1.png}
\end{center}
\caption{\itshape Montaje experimental.}
\label{mifigura1}
\end{figure}

\section*{Resultados}
Los resultados deben presentarse preferiblemente en forma de gráficos y tablas. Además,
es aconsejable expresar resultados con sus incertidumbres\\

Las figuras deben tener buena resolución (mínimo 300 pixeles por pulgada) en escala de grises o a color de ser estrictamente necesario.  Al graficar diferentes series de datos usar en lo posible diferentes símbolos (marcadores), como se muestra en la fig. \ref{mifigura2}.

\begin{figure}[H]
\begin{center}
\includegraphics[width=7cm]{figura2.jpg}
\end{center}
\caption{\itshape Ejemplo figura.}  
\label{mifigura2}
\end{figure}

\begin{table}[H]
\centering
\caption{Ejemplo de tabla. Describa muy bien la tabla, si utiliza símbolos matemáticos aclare su significado. No olvide agregar las unidades de cada variable. Ejemplo: Distancia ($x$) y tiempo para el lanzamiento de cada objeto de prueba caracterizado por su material.}
\label{mitabla}
\begin{tabular}{|l|l|l|}
\hline
Material& Tiempo $\pm 0.1$ $(s)$& $x \pm 1 \times 10^{-2} $ $(m)$ \\ \hline
Aluminio                      & 10.1               & 8                                                 \\ \hline
Bronce                        & 7.0                & 6                                                 \\ \hline
Hierro                        & 6.2                & 9                                                 \\ \hline
\end{tabular}
\end{table}

\section*{Discusión de resultados}
La discusión de datos en artículos científicos corresponde
a la forma como los resultados son interpretados por el
investigador, tanto a la luz de la hipótesis planteada, como
a la de lo que otros autores dicen o han encontrado sobre
el tema. Se trata de dilucidar qué significan los resultados
y por qué ocurrieron de ese modo las cosas. Si se trata de un trabajo experimental, se deben analizar
las dependencias observadas entre las variables y la comparación de los datos con un modelo
propuesto. Si el trabajo además propone un modelo que trate de dar cuenta
de los datos obtenidos, o bien, si se usó un modelo tomado de otros trabajos,
en tal caso debe citarse la fuente consultada.


\subsection*{Conclusiones}
Las conclusiones son las contribuciones del trabajo, de cada discusión de resultados
de la sección anterior debería salir 1 conclusión en lo posible. No escriba conclusiones
que no tenga soporte en sus resultados. Puede listarlas de la siguiente manera:
\begin{itemize}
\item Conclusión 1.
\item Conclusión 2.
\end{itemize}

%- Si se requieren agradecimientos, reconocimientos a entidades, permisos de publicación, etc., irán al terminar el texto y antes de las Referencias bibliográficas.

%archivo tipo bib donde se encuentran las referencias
\bibliography{bibfile}

\end{document}



